% Part: first-order-logic
% Chapter: completeness
% Section: compactness

\documentclass[../../../include/open-logic-section]{subfiles}

\begin{document}

\iftag{FOL}
      {\olfileid{fol}{com}{com}}
      {\olfileid{pl}{com}{com}}

\olsection{સઘનતા પ્રમેય}

પૂર્ણતા પ્રમેયનું એક મહત્વનું પરિણામ સઘનતા પ્રમેય છે. સઘનતા પ્રમેય
કહે છે કે જો !!{sentence}sના ગણનો દરેક \emph{સાન્ત} ઉપગણ સંતોષ્ય હોય,
તો આખો ગણ સંતોષ્ય છે—ભલે ગણ પોતે અનંત હોય. આ જરાય સ્પષ્ટ નથી. પહેલી
નજરે તો એવું કશું દેખાતું નથી જે એવી સંભાવના નકારે કે !!{sentence}sના
કેટલાક અનંત ગણ વિરોધાભાસી હોય, પરંતુ વિરોધાભાસ જાણે અનંત સંખ્યાને
કારણે જ ઊભો થતો હોય. સઘનતા પ્રમેય કહે છે કે આવી પરિસ્થિતિ થઈ શકતી
નથી: !!{sentence}sનો એવો કોઈ અસંતોષ્ય અનંત ગણ નથી જેનો દરેક સાન્ત ઉપગણ સંતોષ્ય
હોય. પૂર્ણતા પ્રમેયની જેમ તેનું ફલિતતા સાથે જોડાયેલું નિરૂપણ પણ છે:
જો !!{sentence}sનો અનંત ગણ કંઈક ફલિત કરે, તો તેનો કોઈ સાન્ત ઉપગણ તે
પહેલેથી જ ફલિત કરે છે.

\begin{defn}
  !!{formula}sનો ગણ~$\Gamma$ \emph{સાન્ત રીતે સંતોષ્ય} છે જો અને માત્ર
  જો દરેક સાન્ત $\Gamma_0\subseteq\Gamma$ સંતોષ્ય છે.
\end{defn}

\begin{thm}[સઘનતા પ્રમેય]
\ollabel{thm:compactness}
!!{sentence}sના કોઈ પણ ગણ~$\Gamma$ અને વાક્ય~$!A$ માટે નીચેના દાવા
સાચા છે:
\begin{enumerate}
  \item $\Gamma\Entails !A$ જો અને માત્ર જો એવો સાન્ત
    $\Gamma_0\subseteq\Gamma$ છે કે $\Gamma_0\Entails !A$.
  \item $\Gamma$ સંતોષ્ય છે જો અને માત્ર જો તે સાન્ત રીતે સંતોષ્ય છે.
\end{enumerate}
\end{thm}

\sourcecorrection{OLCO-014}{મૂળની \texttt{compactness.tex}, પંક્તિ~35માં
``any sentences $\Gamma$ and $!A$'' લખ્યું છે, જોકે~$\Gamma$ વાક્ય
નહિ પણ વાક્યોનો ગણ છે. અહીં ચલોનો પ્રકાર સ્પષ્ટ કર્યો છે.}

\begin{proof}
આપણે~(2) સાબિત કરીએ. જો~$\Gamma$ સંતોષ્ય હોય, તો એવી
\iftag{FOL}{!!a{structure}~$\Struct{M}$}{!!a{valuation}~$\pAssign{v}$}
છે કે દરેક $!A\in\Gamma$ માટે
$\iftag{FOL}{\Sat{M}}{\pSat{v}}{!A}$. અલબત્ત, આ
$\iftag{FOL}{\Struct{M}}{\pAssign{v}}$~$\Gamma$ના દરેક સાન્ત ઉપગણને
પણ સંતોષે છે, તેથી~$\Gamma$ સાન્ત રીતે સંતોષ્ય છે.

હવે ધારો કે~$\Gamma$ સાન્ત રીતે સંતોષ્ય છે. તો દરેક સાન્ત ઉપગણ
$\Gamma_0\subseteq\Gamma$ સંતોષ્ય છે. યથાર્થતા મુજબ
(\iftag{FOL}{%
  \tagrefs{prfAX/{fol:axd:sou:cor:consistency-soundness},
    prfSC/{fol:seq:sou:cor:consistency-soundness},
    prfND/{fol:ntd:sou:cor:consistency-soundness},
    prfTab/{fol:tab:sou:cor:consistency-soundness}}}{%
  \tagrefs{prfAX/{pl:axd:sou:cor:consistency-soundness},
    prfSC/{pl:seq:sou:cor:consistency-soundness},
    prfND/{pl:ntd:sou:cor:consistency-soundness},
    prfTab/{pl:tab:sou:cor:consistency-soundness}}}),
દરેક સાન્ત ઉપગણ સુસંગત છે. હવે
\iftag{FOL}{%
  \tagrefs{prfAX/{fol:axd:ptn:prop:proves-compact},
    prfSC/{fol:seq:ptn:prop:proves-compact},
    prfND/{fol:ntd:ptn:prop:proves-compact},
    prfTab/{fol:tab:ptn:prop:proves-compact}}}{%
  \tagrefs{prfAX/{pl:axd:ptn:prop:proves-compact},
    prfSC/{pl:seq:ptn:prop:proves-compact},
    prfND/{pl:ntd:ptn:prop:proves-compact},
    prfTab/{pl:tab:ptn:prop:proves-compact}}}
મુજબ~$\Gamma$ પોતે સુસંગત હોવો જ જોઈએ. પૂર્ણતા મુજબ
(\olref[cth]{thm:completeness}), ~ $\Gamma$ સુસંગત હોવાથી તે સંતોષ્ય છે.
\end{proof}

\tagprob{FOL}
\begin{prob}
\olref[fol][com][com]{thm:compactness}નો~(1) સાબિત કરો.
\end{prob}
\tagendprob

\tagprob{notFOL}
\begin{prob}
\olref[pl][com][com]{thm:compactness}નો~(1) સાબિત કરો.
\end{prob}
\tagendprob

\iftag{FOL}{%
\begin{ex}
સિદ્ધાંત~$\Gamma$ના દરેક નિદર્શ~$\Struct{M}$માં દરેક પદ~$t$ અલબત્ત
~$\Domain{M}$નો કોઈ !!a{element} નિર્દિષ્ટ કરે છે. શું આપણે એ પણ
ખાતરી આપી શકીએ કે~$\Domain{M}$નો દરેક !!{element} કોઈ ને કોઈ પદથી
નિર્દિષ્ટ થાય? બીજા શબ્દોમાં, શું એવા સિદ્ધાંતો~$\Gamma$ છે કે જેમના
બધા નિદર્શો આવૃત હોય? સઘનતા પ્રમેય બતાવે છે કે જો~$\Gamma$ના અનંત
નિદર્શો હોય, તો એવું નથી. આ જોવા માટે: ધારો કે~$\Struct{M}$ એ
~$\Gamma$નો અનંત નિદર્શ છે અને~$c$ એવું !!a{constant} છે જે
~$\Gamma$ની ભાષામાં નથી. ~ $\Gamma$ની ભાષા~$\Lang{L}$ના દરેક પદ~$t$
માટે વાક્યો~$\eq/[c][t]$નો ગણ~$\Delta$ લો, એટલે કે
  \[
  \Delta = \Setabs{\eq/[c][t]}{t \in \Trm[L]}.
  \]
~$\Gamma\cup\Delta$નો કોઈ સાન્ત ઉપગણ $\Gamma'\cup\Delta'$ સ્વરૂપે
લખી શકાય, જ્યાં $\Gamma'\subseteq\Gamma$ અને
$\Delta'\subseteq\Delta$. ~ $\Delta'$ સાન્ત હોવાથી તેમાં સાન્ત
સંખ્યામાં જ પદો આવી શકે. એવો $a \in \Domain{M}$ લો જે
~$\Domain{M}$નો !!a{element} છે અને તેમાંથી કોઈ પદથી નિર્દિષ્ટ થતો ન
હોય; અને એવી !!{structure}
~$\Struct{M'}$ લો જે~$\Struct{M}$ જેવી જ છે, પણ જેમાં વધુમાં
$\Assign{c}{M'}=a$. ~ $\Delta'$માં આવતા દરેક~$t$ માટે
$a\neq\Value{t}{M}$ હોવાથી $\Sat{M'}{\Delta'}$. વળી
$\Sat{M}{\Gamma}$, $\Gamma'\subseteq\Gamma$, અને~$c$ એ~$\Gamma$માં
ન આવતું હોવાથી $\Sat{M'}{\Gamma'}$. તેથી
$\Gamma\cup\Delta$ના દરેક સાન્ત ઉપગણ $\Gamma'\cup\Delta'$ માટે
$\Sat{M'}{\Gamma'\cup\Delta'}$. આમ~$\Gamma\cup\Delta$નો દરેક સાન્ત
ઉપગણ સંતોષ્ય છે. સઘનતા મુજબ~$\Gamma\cup\Delta$ પોતે સંતોષ્ય છે.
તેથી એવી !!{structure}s~$\Struct{N}$ છે કે
$\Sat{N}{\Gamma\cup\Delta}$. આવી દરેક~$\Struct{N}$ એ~$\Gamma$નો
નિદર્શ છે, પરંતુ આવૃત નથી, કારણ કે~$\Lang{L}$ના દરેક પદ~$t$ માટે
$\Value{c}{N}\neq\Value{t}{N}$.
\end{ex}

\sourcecorrection{OLCO-015}{મૂળની \texttt{compactness.tex}, પંક્તિ~120માં
``there are models $\Sat{M}{\Gamma\cup\Delta}$''થી સંરચનાને સંતોષના
સૂત્ર સાથે વ્યાકરણવિરુદ્ધ રીતે સરખાવવામાં આવી છે અને અગાઉના~$M$નો
પુનઃઉપયોગ થાય છે. અહીં નવી સંરચના~$N$ માટે સંતોષ-દાવો સ્પષ્ટ કર્યો છે.}

\begin{ex}
એવી !!a{language}~$\Lang{L}$ વિચારો જેમાં !!{predicate}~$<$,
!!{constant}s~$\Obj{0}$, $\Obj{1}$ અને !!{function}s $+$, $\times$ તથા
$-$ છે. આ !!{language}માંનાં એવા બધાં !!{sentence}sનો ગણ~$\Gamma$ લો
જે વ્યાપકક્ષેત્ર~$\Rat$ અને સ્વાભાવિક અર્થઘટનોવાળી
!!{structure}~$\Struct{Q}$માં સાચાં છે. આમ~$\Gamma$ એ~$\Lang{L}$નાં
સંમેય સંખ્યાઓ અંગે સાચાં બધાં !!{sentence}sનો ગણ છે. અલબત્ત,~$\Rat$માં
(અને~$\Real$માં પણ) એવી કોઈ સંખ્યા~$r$ નથી જે~$0$ કરતાં મોટી, પરંતુ
દરેક $k\in\PosInt$ માટે $1/k$ કરતાં નાની હોય. જો આવી સંખ્યા હોત, તો તે
\emph{અનંતસૂક્ષ્મ} હોત: શૂન્યેતર, પરંતુ અનંત રીતે નાની. સઘનતા પ્રમેય
વાપરીને બતાવી શકાય છે કે~$\Gamma$ના એવા નિદર્શો છે જેમાં અનંતસૂક્ષ્મો
અસ્તિત્વ ધરાવે છે. આપણી ભાષામાં ભાગાકાર માટે !!a{function} નથી
(શૂન્ય વડે ભાગાકાર અવ્યાખ્યાયિત છે, જ્યારે !!{function}sનું અર્થઘટન
સર્વત્ર વ્યાખ્યાયિત વિધેયોથી કરવું પડે છે). તેમ છતાં $r<1/k$ વ્યક્ત
કરી શકાય છે, કારણ કે આ સાચું હોય જો અને માત્ર જો $r\cdot k<1$. હવે
~$c$ કોઈ નવું !!{constant} અને~$\Delta$ નીચેનો ગણ લો:
\[
\{\Obj{0}<c\}
\cup \Setabs{c\times\num k<\Obj{1}}{k\in\PosInt}
\]
(જ્યાં $\num{k}=(\Obj{1}+(\Obj{1}+\dots+(\Obj{1}+\Obj{1})\dots))$માં
$k$ વાર~$\Obj{1}$ આવે છે). ~ $\Delta$ના કોઈ પણ સાન્ત ઉપગણ
~$\Delta_0$ માટે એવું~$K$ હોય છે કે~$\Delta_0$માં આવતાં બધા
!!{sentence}s $c\times\num{k}<\Obj{1}$ માટે $k<K$. જો~$\Struct{Q}$ને
$\Assign{c}{Q'}=1/K$ સાથે~$\Struct{Q'}$ સુધી વિસ્તારીએ, તો કોઈ પણ
સાન્ત $\Gamma_0\subseteq\Gamma$ માટે
$\Sat{Q'}{\Gamma_0\cup\Delta_0}$; અને તેથી~$\Gamma\cup\Delta$ સાન્ત
રીતે સંતોષ્ય છે (પ્રશ્ન: આ વિગતવાર સાબિત કરો). સઘનતા મુજબ
~$\Gamma\cup\Delta$ સંતોષ્ય છે. ~ $\Gamma\cup\Delta$ના કોઈ પણ નિદર્શ
~$\Struct{S}$માં અનંતસૂક્ષ્મ છે, એટલે કે~$\Assign{c}{S}$.
\end{ex}
}{}

\sourcecorrection{OLCO-016}{મૂળની \texttt{compactness.tex},
પંક્તિઓ~157--158માં ``for all the sentences ... have'' એવી ખામીયુક્ત
રચના છે. અહીં પરિમાણનો વ્યાપ સ્પષ્ટ કરીને દરેક સંબંધિત વાક્યનો સૂચક
$k<K$ છે તેમ લખ્યું છે.}

\tagprob{FOL}
\begin{prob}
અંકગણિતના પ્રમાણભૂત નિદર્શ~$\Struct{N}$માં એવો કોઈ
!!{element}~$k\in\Domain{N}$ નથી જે દરેક સૂત્ર~$\num{n}<x$ને સંતોષે
(જ્યાં~$\num{n}$ એ~$n$~$\prime$ ધરાવતું
$\Obj{0}^{\prime\dots\prime}$ છે). સઘનતા પ્રમેય વાપરીને બતાવો કે
અંકગણિતની ભાષાનાં પ્રમાણભૂત નિદર્શ~$\Struct{N}$માં સાચાં વાક્યો કોઈ
એવી !!a{structure}~$\Struct{N'}$માં પણ સાચાં છે જેમાં દરેક સૂત્ર
$\num{n}<x$ને સંતોષતો !!a{element} \emph{હોય છે}.
\end{prob}
\tagendprob

\iftag{FOL}{%
\begin{ex}
આપણે જાણીએ છીએ કે !!{identity} ધરાવતું પ્રથમ-ક્રમ તર્કશાસ્ત્ર
!!{domain}નું કદ ઓછામાં ઓછું અમુક કદનું હોવું જોઈએ તે વ્યક્ત કરી શકે
છે: વાક્ય~$!A_{\ge n}$ (જે કહે છે કે ``ઓછામાં ઓછા~$n$ જુદા પદાર્થો
છે'') ફક્ત એવી સંરચનાઓમાં સાચું છે જેમાં~$\Domain{M}$માં ઓછામાં ઓછા
~$n$ પદાર્થો હોય. તેથી જો લઈએ
  \[
  \Delta = \Setabs{!A_{\ge n}}{n \ge 1}
  \]
તો~$\Delta$નો દરેક નિદર્શ અનંત હોવો જ જોઈએ. આમ કોઈ સિદ્ધાંતમાં
~$\Delta$ ઉમેરીને ખાતરી કરી શકાય છે કે તેને માત્ર અનંત નિદર્શો જ હોય:
~$\Gamma\cup\Delta$ના નિદર્શો બરાબર~$\Gamma$ના અનંત નિદર્શો છે.

આમ પ્રથમ-ક્રમ તર્કશાસ્ત્ર અનંતતા વ્યક્ત કરી શકે છે. પરંતુ સઘનતા પ્રમેય
બતાવે છે કે તે સાન્તતા વ્યક્ત કરી શકતું નથી. પ્રતિપક્ષે ધારો કે
વાક્યોનો કોઈ ગણ~$\Lambda$ બધી અને માત્ર સાન્ત !!{structure}sમાં
સંતોષાતો હોય. તો~$\Delta\cup\Lambda$ સાન્ત રીતે સંતોષ્ય છે. શા માટે?
ધારો કે $\Delta'\cup\Lambda'\subseteq\Delta\cup\Lambda$ સાન્ત છે,
જ્યાં $\Delta'\subseteq\Delta$ અને $\Lambda'\subseteq\Lambda$. જો
~$\Delta'$ ખાલી હોય, તો કોઈ પણ સાન્ત !!{structure} તેને અને~$\Lambda'$ને
સંતોષે છે. અન્યથા, $!A_{\ge n}\in\Delta'$ હોય એમાં સૌથી મોટી સંખ્યા
~$n$ લો. ~ $\Lambda$ બધી સાન્ત સંરચનાઓમાં સંતોષાતો હોવાથી તેને સાન્ત,
પરંતુ~$\ge n$ !!{element}sવાળો નિદર્શ~$\Struct{M}$ છે. પરંતુ
પછી $\Sat{M}{\Delta'\cup\Lambda'}$. સઘનતા મુજબ
~$\Delta\cup\Lambda$નો અનંત નિદર્શ છે, જે~$\Lambda$ માત્ર સાન્ત
!!{structure}sમાં સંતોષાય છે એ પૂર્વધારણાનો વિરોધાભાસ છે.
\end{ex}
}{}

\sourcecorrection{OLCO-017}{મૂળની \texttt{compactness.tex},
પંક્તિઓ~209--211માં સાન્ત~$\Delta'$માંના સૂચકોમાં સૌથી મોટો~$n$ લેવામાં
આવે છે, પરંતુ ખાલી~$\Delta'$ માટે એવો સૂચક નથી. અહીં ખાલી કિસ્સો
અલગથી નોંધ્યો છે.}

\end{document}
